% from ega4-17-fr.tex
\oldpage[IV-4]{91}

bijectif au point $x'$ (Bourbaki, \emph{Alg. comm.}, chap. II, § 3, n\textsuperscript{o} 2, cor. de la prop. 6), donc aussi dans un voisinage de $x'$ (0$_{\mathrm I}$, 5.2.7); cela entraîne évidemment b), ainsi que la dernière assertion de l'énoncé, en vertu de l'exactitude de (17.13.2.2). Inversement, comme $\alpha\otimes1$ se factorise en
\[
g^*(\sh{N}_{Y/X})\otimes k(x')\xrightarrow{\operatorname{gr}_1(g)\otimes1}\sh{N}_{Y'/X'}\otimes k(x')\xrightarrow{\delta\otimes1}\Omega^1_{X'/S}\otimes k(x')
\]
et que $\operatorname{gr}_1(g)$ est surjectif, dire que $\alpha\otimes1$ est injectif entraîne que $\delta\otimes1$ l'est et que $\operatorname{gr}_1(g)\otimes1$ est bijectif. On en conclut (17.12.1) que $Y'$ est lisse sur $S$ au point $x'$, et que $\operatorname{gr}_1(g)$ est bijectif dans un voisinage de $x'$ (Bourbaki, \emph{loc. cit.}). En outre, comme la suite (17.13.2.2) est alors exacte, et que $\sh{N}_{Y'/X'}$, isomorphe à $g^*(\sh{N}_{Y/X})$, est de rang $c$ au point $x'$, $\Omega^1_{Y'/S}$ est de rang $n-c$ en ce point, ce qui achève de prouver l'équivalence de a) et b), en vertu de (17.10.2).

Reste à montrer l'équivalence de b) et b') lorsque $x'$ est rationnel sur $k(s)$ (ce qui implique qu'il en est de même de $x$); alors $g^*(\sh{N}_{Y/X})\otimes_{\sh{O}_{Y'}}k(x')$ s'identifie à $\sh{N}_{Y/X}\otimes_{\sh{O}_Y}k(x')$, et puisque la suite
\[
0\to\sh{N}_{Y/X}\to\Omega^1_{X/S}\otimes_{\sh{O}_X}\sh{O}_Y\to\Omega^1_{Y/S}\to0
\]
est exacte au point $x$ et formée de $\sh{O}_Y$-Modules localement libres, le dual de $\sh{N}_{Y/X}\otimes_{\sh{O}_Y}k(x)$ s'identifie à l'espace quotient $\mathrm{T}_{X/S}(x)/\mathrm{T}_{Y/S}(x)$ (16.5.12); d'où l'équivalence de b) et b').

\begin{definition}[17.13.3]
\EGAReaderTarget{ega4.c04.def.17-13-3}{IV.17.13.3}
Avec les notations de (17.13.1), on dit que le morphisme $f$ est \emph{transversal} à $Y$ au point $x'$, relativement à $S$, si $X$ et $Y$ sont lisses sur $S$ au point $x$, $X'$ lisse sur $S$ au point $x'$, et si les conditions équivalentes a), b) de (17.13.2) sont satisfaites.
\end{definition}

Lorsque aucune confusion n'est à craindre, on supprime la mention du préschéma $S$ et on dit simplement que $f$ est transversal à $Y$ au point $x'$.

\begin{remarks}[17.13.4]
\label{IV.17.13.4}
(i) Supposons que $X$, $Y$ et $X'$ soient \emph{plats et localement de présentation finie} sur $S$. Pour tout $s\in S$, notons $X_s$, $Y_s$, $X'_s$, $Y'_s$ les fibres de $X$, $Y$, $X'$, $Y'$ au point $s$, $f_s:X'_s\to X_s$ le morphisme déduit de $f$ par changement de base. Il résulte alors de (17.5.9) que pour que $f$ soit un morphisme transversal à $Y$ au point $x'$, il faut et il suffit que, si $s$ est l'image de $x$ dans $S$, $f_s$ soit transversal à $Y_s$ au point $x'$, relativement à $k(s)$.

(ii) On a vu dans (17.13.2) que l'ensemble des points $x'\in Y'$ où $f$ est transversal à $Y$ est \emph{ouvert} dans $Y'$; si $Y'$ est de plus \emph{propre} sur $S$, on en déduit que l'ensemble des $s\in S$ tels que $f$ soit transversal à $Y$ en tous les points de $Y'_s$, relativement à $S$, est \emph{ouvert} dans $S$.

Lorsque $X$, $Y$ et $X'$ sont plats et localement de présentation finie sur $S$, il résulte de (i) que l'ensemble des $x'\in Y'$ tels que $f_s$ soit transversal à $Y_s$ au point $x'$ ($s$ image de $x'$ dans $S$), relativement à $k(s)$, est \emph{ouvert} dans $Y'$. Si de plus $Y'$ est propre sur $S$, l'ensemble des $s\in S$ tels que $f_s$ soit transversal à $Y_s$ en tous les points de $Y'_s$ (relativement à $k(s)$) est \emph{ouvert} dans $S$.

(iii) La propriété d'être lisse sur $S$, ainsi que la notion de dimension relative à $S$ en un point, étant stables par tout changement de base $S'\to S$ ((17.3.3) et (4.2.7)), il en est de même de la propriété pour un morphisme $f:X'\to X$ d'être transversal à un sous-préschéma de $X$ en un point, relativement à $S$.

